\section{Higher Order Derivatives \textcolor{red}{(to be deleted)}}
\label{sec:high-order-deriv}
\markright{Derivatives of Derivatives}

We have seen that if, for example, $y=x^3,$ then $\dfdx{y}{x} = 3x^2.$
There are at least three distinct  ways to interpret the formula $\dfdx{y}{x} =
3x^2.$
The first two of these are the ones we've been using all along:
\begin{description}
\item[Geometric Interpretation:] The derivative $\dfdx{y}{x} = 3x^2$ is the formula
  that tells us the slope of the curve it was ``derived\footnote{Hence the
  name.}'' from, $y=x^3.$
\item[Dynamic Interpretation] If $y=x^3$ is some 
  relationship between two quantities (like position and time, for
  example) then $\dfdx{y}{x}$ tells us how fast $y$ is changing
  with respect to $x.$ (Just as velocity is how fast position is changing
  with respect to the elapsed time.)
\item[Functional (Modern) Interpretation] A relation like $y=x^3$ gives $y$
  \underline{as a function of} $x.$ What this means, really, is that if
  a value for $x$ is given then the value of $y$ is forced. For
  example if $x=2$ then $y=2^3 = 8.$ If we choose any other value for
  $x$ then $y$ is forced to be the cube of $x.$ We usually encapsulate
  this idea with the notation $y(x) = x^3,$ so that $y(2)=2^3=8;$
  $y(3)=3^3=27;$ $y(4)=4^3=64,$ and so on.

 It should be clear
  that both of the formulas $y(x) = x^3$ and $\dfdx{y}{x} = 3x^2$
  represent functions in the sense just given: If $x$
  is given then the values of both $y$ and $\dfdx{y}{x}$ are
  forced. For example, if $x=2$ then $y=2^3=8$ and
  $\dfdx{y}{x}= 6(2^2) = 24.$ That is, if $x$ is chosen, then $y$ and
  $\dfdx{y}{x}$ are both forced. 

  This gives us a new way to think about the expression $\dfdx{y}{x}.$
  We can think of it as a function of $x$ just as $y$ is a function of
  $x.$ It is true that $\dfdx{y}{x}$ is \emph{derived}\footnote{Hence
      the name.} from $y$ in the same way and for the same reasons we
  have been considering all along. So the first two interpretations
  are still valid and there is no reason not to use them as
  appropriate. But we will also find this third interpretation to be
  useful as well
\end{description}

\begin{mynotation}{Function}
  Interpreting the symbol $\dfdx{y}{x}$ as the name of a function can
  lead to some confusion so we will take a moment to address this
  before we proceed.

  Until now we have interpreted the symbol $\dfdx{y}{x}$ as the ratio
  of the differentials $\dx{y}$ and $\dx{x}.$ This will remain a very
  useful interpretation and we will continue to use it. But it will
  also be useful to interpret the symbol $\dfdx{y}{x}$ as the name of
  the function derived from $y,$ just as the symbols ``$\sin$'' and
  ``$\cos$'' are the names of certain trigonometric functions. As such
  ``$\dfdx{y}{x}$'' is \alert{one symbol} and cannot be separated into
  the constituent parts $\dx{y}$ and $\dx{x}.$

  It should be clear at this point that the derivative is by far the
  more important concept. This is why we introduced the notion of
  differentials as ``useful fictions'' back in
  Section~\ref{sec:some-rules-diff}. We only really need differentials
  as stepping stones on our way to building the derivative. To be
  sure, we will continue to use them because they are very ``useful
  fictions.'' But henceforth they will only appear as an intermediary
  step in building the derivative.
\end{mynotation}

If the symbol $\dfdx{y}{x}$ is the name of the function we get by
differentiating $y$ then a natural question is, ``What is the
derivative of the function $\dfdx{y}{x}?$'' Before we start down that
path though, let's first make sure we don't get sloppy. If we can take
more than one derivative of $y$ we can't just use the word
``derivative'' to refer all of them. We'll need a finer distinction
than that.


The first time  we differentiate $y(x) = x^3$  we get $\dfdx{y}{x}.$
The second time we differentiate $y(x) = x^3$ we
 get $6x.$ What should we call
$6x?$ The obvious
choice is to call $6x$ the \emph{second derivative} of $y=x^3.$
Clearly $3x^2$  is the \emph{first derivative\footnote{But we will
    continue to call it ``the derivative'' when there is no chance of
    confusion.}} and $6$ is the
\emph{third derivative} of $y=x^3.$

All well and good, but what it will be cumbersome to write out ``the
third derivative'' every time we need it. What \underline{notation}
should we use to indicate higher order derivatives?


For better or worse the extension of the differential\footnote{Several
  notations have been invented. However, we've been using the
  differential notation all along and we see no value in changing the
  notation on you at this point. So we will continue to use the
  differential notation for now.

  % However, mathematical notation is a tool and, like all tools, you'll
  % need to use the one that was designed for the task at hand. By
  % limiting you to only the differential notation we are necessarily
  % limiting the tool set available to you. Once you have fully mastered
  % the \emph{concept} differentiation we will introduce the other
  % notations.
} notation to higher order derivatives that has been
become standard is this:
\begin{align*}
  \text{ If } y &= x^3 \text{ then}\\
  \dfdx{y}{x}&= 3x^2 \text{ as before, and}\\
  \dfdxn{y}{x}{2} &= 6x.\\
\intertext{Differentiating one more time we get}
  \dfdxn{y}{x}{3} &= 6.
\end{align*}

Notice the odd placement of the exponents. This is an historical
accident. Leibniz would have proceeded as follows:
\begin{align}
\nonumber  \text{Let } y&= x^3.\\
  \intertext{Then}
\nonumber  \dx{y}&= \dx{(x^3)}\\
\nonumber               &= 3x^2\dx{x}.\\
  \intertext{So far, so good. But in taking the second differential
  Leibniz would have used the Product Rule, giving}
\nonumber  \dx{(\dx{y})} &= 3x^2\underbrace{\dx{(\dx{x})}}_{=0} +
                                       6x\dx{x}\dx{x}\\
  \intertext{so}
\label{eq:PRNotation} \dx{(\dx{y})} &=    6x\dx{x}\dx{x}.
\end{align}
To see why  $\dx{(\dx{x})}=0$ recall that, whatever else it might
  be, the differential $\dx{x}$ is a constant. And the differential of
  any constant is zero, by our first differentiation rule.

  A natural way to abbreviate the left side of
  equation~\ref{eq:PRNotation} is $\dx{(\dx{y})}=\dx^2y$ and the on the
  right side the expression $\dx{x}\dx{x}$ is just a product so it is equal
  to $\dx{x}^2.$ Thus when we divide through by $\dx{x}^2$ we get
  $\dfdxn{y}{x}{2}=6x.$

  Similarly we get $\dfdxn{y}{x}{4} = 0.$ Obviously we can continue
  this process as many times as we would like, although for this
  particular example every derivative after the fourth is zero so this
  would be rather dull. The following example is more interesting.

\begin{myexample}
  Consider the expression $y=\sin(x).$ As we've seen $\dfdx{y}{x}=
  \cos(x).$ Continuing we have $\dfdxn{y}{x}{2} = -\sin(x),$
  $\dfdxn{y}{x}{3} = -\cos(x),$ and finally $\dfdxn{y}{x}{4} =
  \sin(x).$ Thereafter the pattern of derivatives clearly repeats.
\end{myexample}

\begin{embeddedproblem}{}
  \begin{enumerate}[label={\bf{}(\alph*)}]{}
  \item   Suppose $y=\sin(x).$ Compute
    \begin{enumerate}[label={\bf{}(\roman*)}]{}
      \begin{multicols}{4}
      \item $\displaystyle\dfdx{y}{x}$
      \item $\displaystyle\dfdxn{y}{x}{2}$
        \item $\displaystyle\dfdxn{y}{x}{3}$
        \item     $\displaystyle\dfdxn{y}{x}{4}$
        \item $\displaystyle\dfdxn{y}{x}{10}$
        \item $\displaystyle\dfdxn{y}{x}{24}$
        \item $\displaystyle\dfdxn{y}{x}{101}$
        \end{multicols}
      \end{enumerate}
      \begin{multicols}{2}
      \item Repeat part (a), using $y=\cos(x).$
      \item Repeat part (a), using $y=e^x.$
      \item Repeat part (a), using $y=xe^x.$
      \item Repeat part (a), using $y=\sin(ax),$ where $a$ is a
        constant.
      \end{multicols}
  \end{enumerate}
\end{embeddedproblem}

\begin{mynotation}{Troubling Aspects of Differentials and Differential}
  We would be remiss if we did not point out some very troubling
  aspects of the  notation we've been using.
  Consider, when we introduced the concept of a differential in
  Chapter \ref{chapt:differentials} we appealed to your understanding
  of \emph{finite} differences. That is, if $x_1\lt{}x_2$ then the finite
  difference
$$
\Delta{}x=x_2-x_1
$$
is a measure of the actual separation between $x_2$ and $x_1.$
Extending this to infinitesimals,  we imagined that $x_2$ and $x_1$ are
\underline{infinitely} close together -- but still distinct -- and we
represented the distance between them as
$$
\dx x = x_2-x_1.
$$

It is all well and good to \emph{write down} expressions like $\dx y$
or $\dx^2 y$ but what do they represent, really?

If $x_1$ and $x_2$ are infinitely close together doesn't that mean
that they are \underline{not} separated at all? If they are not
separated then $x_2=x_1.$ In that case isn't $\dx x =x_2-x_1$ just
zero? And if $\dx x =0$ then doesn't the expression $\dfdx{y}{x}$ call
for division by zero? Which we know is meaningless?

Now consider the meaning of the \underline{second} derivative. If the
first derivative is obtained by computing the ratio of quantities that
are infinitely close together, what can it possibly mean to
differentiate it again? If we've already taken numbers that are
infinitesimally close together to compute the first derivative, what
can we possibly be using to compute the second, or the third
derivative?

What is the infinitesimal difference of the infinitesimal difference of
an infinitesimal difference?

This is all very puzzling and it is tempting to toss it all away.
Except that at this point we have a great deal of evidence that
$\dfdx{\left(x^3\right)}{x}=3x^2,$ and
$\dfdx{\left(3x^2\right)}{x}=6x$ are both quite meaningful. Among
other things we've used them to find the equations of tangent lines to
a variety of curves quite successfully. There must be something to
this. But what?

If we are honest with ourselves we must acknowledge the legitimacy of
these questions. There is something here that we do not yet
understand. We will proceed very cautiously.
\end{mynotation}


\begin{Digression}[Approximation with Polynomials]
\begin{wrapfigure}[]{r}{2in}
  \includegraphics*[height=1.5in,width=2in]{../Figures/CosApprox2}
  \label{fig:CosApprox2}
  \captionsetup{labelformat=empty}
\end{wrapfigure}
  When $-\frac{\pi}{2}\le x\le \frac{\pi}{2}$ a well known
  approximation of $\cos(x)$ is
  \[
    \cos(x)\approx 1-\frac{x^2}{2}+\frac{1}{24}x^4.
  \]
  It is clearly a pretty good approximation; when you compare the
  graphs of $\cos(x)$ and $1-\frac{x^2}{2}$ it is clear that they are
  close together, at least for small values of $x.$ But how in the
  world would anyone come up with this approximation in the first
  place?

  This is actually a straightforward generalization of the 
  construction of tangent lines using the derivative.
  We begin by constructing the tangent line of this graph at the point
  \((0,1).\) This is not hard. If \(y=\cos(x)\) then
  \(\left.\dfdx{y}{x}\right|_{x=0}=0. \) Therefore an equation of the
  tangent line at \((0,1)\) is
  \begin{equation}
    y=\left[\eval{\dfdx{y}{x}}{x}{0}\right](x-0)
    +1\label{eq:CosTanAtZero}
  \end{equation}
  or
  \[y=1.\] Most likely you knew that already.

\begin{wrapfigure}[]{r}{2in}
\captionsetup{labelformat=empty}
  \centerline{\includegraphics*[height=1in,width=2in]{../Figures/CosApprox}}
\label{fig:CosApprox}
\end{wrapfigure}
From the figure at the right we see that at the point of tangency,
$(0,1)$, the graphs of $y=\cos(x)$ and $y=1$ are both
horizontal. Indeed, this should be fairly obvious. That's what a
tangent line is, a line with the same slope as the curve it is tangent
to. But suppose we could find a function whose first \emph{and second}
derivatives match $\cos(x)$ at $x=0?$ What would that look like?  That
function is $y=1-\frac{x^2}{2}$ and it looks like the blue graph
below:
\centerline{\includegraphics*[height=2in,width=2in]{../Figures/CosApprox2}}

\begin{myexample}
  You can see where this is going, right? How could we find the polynomial
  whose first four derivatives match the
  $\cos(x)$ at $x=0$? % The curve
  % $y= 1-\frac{x^2}{2}+\frac{x^4}{24}$ is the polynomial with the
  % property that its first four derivatives at $x=0$ match the first
  % four derivatives of $\cos(x),$ also at $x=0.$


  % But we did not ask ``What is the function?'' We asked, ``How do we
  % find it?''
  This is actually pretty straightforward. A general
  quartic\footnote{A quartic is a fourth degree polynomial.} has the
  form:
$$
y(x) = A+Bx+Cx^2+Dx^3+Ex^4.
$$
So all we have to do is determine what the numbers $A,$ $B,$ $C,$ $D,$
and $E$ have to be. We do this in five steps:
\begin{enumerate}
\item First we want $\cos(0)$ to be equal to $y(0).$ Evaluating each
  we see that $\cos(0)=1$ and $y(0)=A.$ Therefore $A=1.$
\item Second we want
  $\left.\left.\dfdx{(\cos(x))}{x}\right|_{x=0}=\dfdx{y}{x}\right|_{x=0}.$
  Evaluating each we have
  $\left.\dfdx{(\cos(x))}{x}\right|_{x=0}=\sin(0)=0,$ and
  $\left.\left.\dfdx{y}{x}\right|_{x=0}=B+2Cx+3Dx^2+4Ex^3\right|_{x=0}=B,$
  so $B=0.$
\item At this point we assume that you've noticed the pattern so we
  will abbreviate our explanations. Setting
  $\left.\left.\dfdxn{(\cos(x))}{x}{2}\right|_{x=0}=\dfdxn{y}{x}{2}\right|_{x=0}$
  gives us
  $\left.\left.-\cos(x)\right|_{x=0}=-1=2C+6Dx+12Ex^2\right|_{x=0}=2C.$
  So, $C=-\frac12.$
\item
  $\left.\left.\dfdxn{(\cos(x))}{x}{3}\right|_{x=0}=\dfdxn{y}{x}{3}\right|_{x=0}$
  gives us
  $\left.\left.-\sin(x)\right|_{x=0}=0=2C+6Dx\right|_{x=0}=2C,$ so
  $C=0.$
\item In the next step we differentiate once more, set the fourth
  derivatives equal, and find that $E=\frac{1}{24}.$
\end{enumerate}
Putting this all together we have
$$
\cos(x)\approx 1-\frac12x+\frac{1}{24}x^4
$$
for values of $x$ near zero.
\end{myexample}

\begin{embeddedproblem}{}
  \begin{enumerate}[label={\bf{}(\alph*)}]
  \item Show that if we require the first six derivatives to match at
    $x=0$ the polynomial we get is
    $y(x) = 1-\frac12x^2+\frac{1}{24}x^4-\frac{1}{720}x^6.$
  \item What happens if we require the first seven derivatives to
    match?
  \end{enumerate}
\end{embeddedproblem}

\begin{embeddedproblem}{}
  \begin{enumerate}[label={\bf{}(\alph*)}]
  \item Find the cubic polynomial whose first two derivatives match
    the first two derivatives of $\sin(x)$ at $x=0.$
  \item Now find the quintic\footnote{Fifth degree polynomial.} whose
    first four derivatives match.
  \item  Repeat parts (a) and (b) but this time make the derivatives
    match at $x=\frac{\pi}{2}.$
  \end{enumerate}
\end{embeddedproblem}

\begin{embeddedproblem}{}
  \begin{enumerate}[label={\bf{}(\alph*)}]
  \item Find the cubic polynomial whose first two derivatives match
    the first two derivatives of $e^x$ at $x=0.$
  \item Now find the quintic whose
    first four derivatives match.
  \item  Repeat parts (a) and (b) but this time make the derivatives
    match at $x=\ln(2).$
  \end{enumerate}
\end{embeddedproblem}
\end{Digression}

\begin{mynotation}{Prime}
  You surely noticed how awkward the notation was in
  equation~(\ref{eq:CosTanAtZero}). The problem is here:
  $\eval{\dfdx{y}{x}}{x}{0}.$ This expression means, ``the value
  of the derivative of $y$ at the point $x=0.$'' To be clear, there is
  nothing wrong with the notation itself. It's just that we will need
  to refer to the value of a derivative at specific points a lot, and
  this notation is just too bulky. We need something sleeker; not so
  hard to write down, and not so hard to read.
  
  As we noted previously, if $y$ is a function of $x$ then
  $\dfdx{y}{x}$ is also a function of $x.$ Specifically, it is the
  function that tells us the slope of the graph of $y(x)$ (geometric
  viewpoint), or the rate of change of $y$ with respect to $x$
  (dynamic viewpoint). We usually indicate the functional nature of
  $y$ by appending ``$(x)$'' to the right, as: $y(x).$ So
  $\eval{\dfdx{y}{x}}{x}{0}$ can be expressed as $\dfdx{y}{x}(0).$
  But this really isn't much of an improvement, is it?

  We've been
  using the differential ratio, $\dfdx{y}{x}$  to denote the
  derivative of $y$ because we originally viewed
  $\dfdx{y}{x}$ as the \emph{actual} ratio of the differentials
  $\dx{y}$ and $\dx{x}.$ But we've moved on from that. We're now
  thinking of the expression $\dfdx{y}{x}$ as just the name of a
  particular function. But names don't have to be so complicated.

  Joseph-Louis Lagrange $(1736-1813)$ introduced the notation,
  $y^\prime(x) = \dfdx{y}{x}.$ (Pronounced ``$y$ prime of $x$.'')
  This notation fills our need very handily. Since
  $ \eval{\dfdx{y}{x}}{x}{0} = y^\prime(0)$ equation(~\ref{eq:CosTanAtZero}) can
  be restated as
  $$
  y=y^\prime(0)(x-0)
  $$
  which is certainly more compact, though it is probably harder to
  read at first. This will get better with practice.

  Bringing the ``prime'' notation into play also begins to move us
  into a more modern way of thinking about Calculus and, in
  particular, the derivative. Previously we emphasized the roles of
  the slope of a graph, and (equivalently) the rate of change of
  relative quantities. The differential quotient $\dfdx{y}{x}$ was just
  the expression that provided that slope (or rate of change).

  To be sure, the derivative still does that and this is
  precisely what makes it important. However it is very difficult to
  understand what the  second, or third, or higher order derivatives
  might be if we continue to think in terms of the differentials
  $\dx{y}$ and $\dx{x}.$ If $\dx{y}$ is an infinitely small difference
  in the quantity $y$ (think of $y$ as a position), then what in the
  world is $\dx(\dx{y})=\dx^2{y}?$ An infinitely small difference
  between infinitely small differences? That sounds a little crazy,
  doesn't it?

  Now ask yourself what $\dx^5y$ might be? Or $\dx^{100}y?$

  On the other hand if $y(x) = x^3,$ then we can ``derive'' the
  derivative $y^\prime(x)=3x^2$ and there are no troublesome
  differentials in the way. We've simply started with one function:
  $y(x) = x^3,$ and, via the rules we've already learned, derived
  another: $y^\prime(x) = 3x^2.$ Since $y^\prime$ is just another
  function there is no logical objection to applying the rules again
  to get $y^{\prime\prime}(x) = 6x.$ Trying to justify this second
  differentiation using differentials is very difficult.

  Finally, it should be clear that Lagrange called $y^\prime$ ``the
  derivative'' because he ``derived'' it from $y.$ So the
  conventions of the English language would dictate that, starting
  with $y(x)=x^3$ we should derive the derivative, $y^\prime(x)=3x^2.$

  Unfortunately the history of our topic interferes with this.  For
  the first $100$ years or so ``the derivative'' didn't exist. In
  those years what came to be called the derivative was simply the
  differential ratio, $\dfdx{y}{x}.$ So we did not ``derive'' to obtain
  a derivative in those years, we ``differentiated'' to obtain
  differentials. After Lagrange, the concept of the derivative as a function
  began to replace the older differential ratio concept. But the
  older way of expressing the computation remained in place. So the modern
  formalism is this: We ``differentiate'' the function $y=x^3$ to
  obtain the ``derivative'' $y^\prime(x)=3x^2.$

  As English this doesn't really make sense, but your teacher will cringe and
  give you the stinky-eye if you don't say it this way.

  Really! Try it and see.

  We will continue to use Leibniz' differential notation in most
  cases. But when it becomes cumbersome (as it will in
  section~\ref{sec:newtons-method-1}) we will switch to Lagrange's
  notation.
\end{mynotation}

% \section{How to Make Things ``Easier on  the Eyes'' (Also Known as the Chain Rule)}
% \label{sec:diff-with-subst}
% \markright{Chain Rule}

% Recall that Fermat was able to deduce Snell's Law of Refraction by
% minimizing the amount of time it took for light to travel from Point
% $A$ in air to point $B$ in water.  Fermat was able to do this before
% the invention of of Calculus because he was a \emph{very} clever man.
% Leibniz reduced it to a fairly routine computation by applying his
% differential Calculus.  Recall Problem~\ref{PIC:FermatSnell} on
% page~\pageref{PIC:FermatSnell} from
% Chapter~\ref{chapter:ScienceBeforeCalc}.

%   Consider the following labeling in the above diagram.
  
% \centerline{\includegraphics*[height=2in,width=4in]{../Figures/Refraction3}}

% Find an expression for the time $T$ traveled along the path $A$ from
% to $B$ in terms of the variable and the constants .

% It turns out that the expression is given by
% $$
% T=T(x)=\frac{\sqrt{a^2+x^2}}{v_1} +\frac{\sqrt{b^2+(c-x)^2}}{v_2}.
% $$
% As with Fermat, Leibniz knew that if we graphed this curve, then at
% the minimum, there would be a horizontal tangent line.  This says that to find this minimum
% we should set $\dfdx{T}{x}=0.$  Applying our rules for
% differentiation, we get
% $$
% \dx{T}=\frac{1}{v_1}  \dx((a^2+x^2 )^{\frac{1}{2}} )+\frac{1}{v_2}
% \dx((b^2+(c-x)^2 )^{\frac{1}{2}})
% $$

% We seem to have run into an impasse. We don't have a rule for
% differentiating $(a^2+x^2)^{\frac12}$.
% Or do we?

% Suppose we make the substitution $z=a^2+x^2$.  Then we have,
% \begin{align*}
% \dx\left((a^2+x^2)^{\frac12}\right)&=\dx(z^{\frac12})=\frac12\left(a^2+x^2\right)^{\frac12}\dx(a^2+x^2) \\
%   &=\frac12\left(a^2+x^2\right)^{\frac12}(2x)\dx  x \\
%   &=\frac{x\dx x}{\sqrt{a^2+x^2}}
% \end{align*}

% Applying the same idea we let $w=b^2+(c-x)^{2}$ and obtain
% \begin{align*}
%   \dx
%   \left(
%   (b^2+(c-x)^2)^{\frac12}
%   \right) &= \dx\left(w^{\frac12}\right) \\
%   &= \frac12w^{\frac12}\dx w \\
%   &= \frac12  (b^2+(c-x)^2)^{\frac12}\dx(b^2+(x-x)^2).
% \end{align*}
% \begin{embeddedproblem}{}
%   \begin{enumerate}[label={\bf{}(\alph*)}]{}
%   \item Finish the  computation of
%     $\dx\left(b^2+(c-x)^2\right)^{\frac12}.$
%   \item Compute $\dfdx{T}{x}$ and set this equal to zero to derive
%     Snell's Law:
%     $$
%     \frac{\sin\theta_1}{v_1}=\frac{\sin\theta_2}{v_2}
%     $$
%   \end{enumerate}
% \end{embeddedproblem}

% The substitution we used above is so useful that it has its own name,
% the Chain Rule.  Many Calculus teachers and students associate this
% name totally with a rule about differentiation.  Surprisingly, the
% name is older than that, but does have relevance to our purposes in
% calculus

% Before the invention of Calculus arithmetic primers gave the following 
% simple rule to be used to, among other things, convert money from one
% currency to another: Suppose we have $x$ dollars  that we want to
% convert to British pounds. If we know that 
% \begin{equation}
%  z \text{ dollars}  = y \text{ pounds}
%  \label{eq:Dollar2PoundConversion}
%  \end{equation}
%  then if we divide both sides of equation(~\ref{eq:Dollar2PoundConversion}) by ($z$ dollar) we get 
% $$
% 1=\frac{y \text{ pounds}}{z \text{ dollars} }.
% $$
% Notice that the left side of this last formula is just the number one. So we start with $x$ dollars, multiply by one in the form $1=\frac{y \text{ pounds}}{ z \text{ dollars} }$ to get
% \begin{align*}
%   x \text{ dollars} &= x\text{ dollars}\cdot1 = x\text{ dollars}\cdot\frac{y \text{ pounds}}{ z \text{ dollars} }\\
%     &= x\cdot y/z\cdot\frac{\text{\cancel{dollars}} \cdot\text{ pounds}}{\text{\cancel{dollars}}}
% %   xy \text{ pounds} &= \left(x \text{ dollars}\right)\left(
% %                     \frac{y \text{ pounds}}{\text{dollars}}\right)\\
% %   &= xy\ (\text{dollars})\left(\frac{\text{pounds}}{\text{dollars}}\right)\\
% \intertext{and since the dollars cancel (at least symbolically) we have}
%   &= xy/z\text{ pounds.}
% \end{align*}

% Now  suppose that $1 \text{ dollar}=.85\text{ pounds},$ and that $1 \text{ pound} = 221 \text{ Chinese yuan}.$ Then\footnote{We're making these numbers up. Don't use them to convert
% currency.} if we
% need to convert $25$ American dollars to Chinese yuan we compute as follows:
% \begin{align*}
%   25\text{ dollars} &= 25\text{ dollars}\times \frac{.85\text{ 
%                       pounds}}{1 \text{ dollar}}\times \frac{221
%                       \text{ 
%                       yuan}}{1 \text{ pound}}\\
%                     &= 25\times.85\times221\left(\text{\cancel{dollars}}\right)
%                       \left(\frac{\text{\bcancel{pounds}}}{\text{\cancel{dollars}}}\right)
%                       \left(\frac{\text{yuan}}{\text{\bcancel{pounds}}}\right)  \\
%                     &= 4696.25 \text{ yuan}  
% \end{align*}

% The\footnote{This table does not belong here. \index{Differentiation Rules!Inverse Trig Functions}
% \begin{center}
%   \begin{tabular}{ |c|c||c|c| }
%     \hline
%     \multicolumn{4}{|c|}{} \\
%     \multicolumn{4}{|c|}{\Large\bf{}Differentiation of } \\
%     \multicolumn{4}{|c|}{} \\
%     \multicolumn{2}{|c||}{\Large\bf{}Trig Functions}  & \multicolumn{2}{|c|}{\Large\bf{} Inverse Trig Functions} \\
%     \multicolumn{2}{|c||}{} &    \multicolumn{2}{|c|}{} \\
%     \hline
%     $y$   &$\dfdx{(y)}{x}$      &$y$&$\dfdx{(y)}{x}$\\
%     \hline\hline
%     $\sin x$ &$ \cos x $        &$\inverse\sin x$ &
%                                                     $\frac{1}{\sqrt{1-x^2}} $\\\hline
% $\cos x$&$ -\sin x $& $\inverse\cos x$ & $\frac{-1}{\sqrt{1-x^2}}$ \\\hline
% $\tan x$  & $\sec^2 x$      &$\inverse\tan x$&$o\frac{1}{1+ x^2}$\\\hline
% $\cot x$  &$-\csc^2 x$      &$\inverse\cot x$&$\frac{-1}{1+
%                                                x^2}$\\\hline
% $\sec x$ &$ \sec x\tan x $&$\inverse\sec x$&$\frac{1}{\abs{x}\sqrt{x^2-1}}$\\\hline
% $\csc x$  &$-\csc x\cot x $ &$\inverse\csc x$&$\frac{-1}{\abs{x}\sqrt{x^2-1}}$\\\hline
%   \end{tabular}
% \end{center}
% } crucial observation here is that in the expression
% $\left(\text{dollars}\right)
% \left(\frac{\text{pounds}}{\text{dollars}}\right)
% \left(\frac{\text{yuan}}{\text{pounds}}\right) $ we have this symbolic
% cancellation of both dollars and pounds so that when finished only
% yuan are left.  This daisy chain of symbolic cancellation was called
% the Chain Rule and it was very helpful in a time when currency
% exchange was done by people. Nowadays of course,  these computations
% are done by computers so a name is  no longer needed. Machines don't
% care about names.

% However, and especially when the differential notation is used to
% express the derivative, the same sort of symbolic cancellation is seen
% to occur.

% We are accustomed to ``cancelling out'' $\dfdx{y}{x}\cdot\dfdx{x}{t}$
% to get $\dfdx{y}{t}.$ Here we want to ``uncancel'' the $\dx{x}$ to get
% $$
% \dfdx{y}{t} = \dfdx{y}{x}\dfdx{x}{t}.
% $$

% More formally, if $y$ depends on $x$ which depends on
% $t$ and we wish to compute $\dfdx{y}{t}$ we proceed as follows:
% \begin{align*}
%   \dfdx{y}{x}        &=   \dfdx{y}{x}\\
%         \dx y  &= \dfdx{y}{x}\dx x\\
% \intertext{and dividing by $\dx t$ on both sides gives,}
%   \dfdx{y}{t} &= \dfdx{y}{x}\dfdx{x}{t}.\\
% \end{align*}

% Notice that \emph{symbollically} we don't seem to have accomplished
% much. On the right side of this last
% formula the $\dx x$ that appears in the numerator appears to cancel out
% the one in the denominator and in the end we have the (apparently)
% obvious statement that $\dfdx{y}{t}= \frac{\dx
%   y}{\text{\cancel{$\dx x$}}}\frac{\text{\cancel{$\dx x$}}}{\dx t} =
% \dfdx{y}{t}.$

% However these symbols we're using are much like the pebbles used by
% medieval merchants. At one level they are just empty symbols (pebbles)
% lying on the page (counting board) and it is
% profoundly obvious that $\dfdx{y}{t} = \dfdx{y}{x}\dfdx{x}{t}.$ Just
% do the cancellation. Duh! 

% But on a counting board a pebble lying here means ``one'' and the same
% pebble lying over there means ``ten.'' To calculate on a counting
% board we have to interpret the meaning of each pebble based on its
% location on the board. 

% We have to interpret the meaning of our Calculus symbols
% similarly. Thus at one level the un-cancellation,
% $\dfdx{y}{t} = \dfdx{y}{x}\dfdx{x}{t}$ is very simple to understand,
% at another, more abstract, level it has a much deeper significance.

% We understand the symbol $\dfdx{y}{t}$ to mean ``the rate of change of
% $y$ with respect to $t$'' so this formula, when properly interpreted
% says that the rate of change of $y$ with respect to $t$ is equal to
% the rate of change of $y$ \underline{with respect to $x$} multiplied
% by the rate of change of $x$ \underline{with respect to $t.$}


% For example\footnote{Somewhere in the Chain Rule section there should
%   be a discussion of ``Substitution'' along the following lines:

%   It seems that we've reached an impasse.  We don't yet know how to
% differentiate the complex expression involved.  In addressing this, we
% will uncover a surprisingly simple, yet powerful tool in all of
% mathematics: substitution.

% The idea behind a substitution is so elementary that it is often taken
% for granted, but we will demonstrate its power here.  Performing a
% substitution does not ``solve'' a problem, but merely takes an
% existing problem and makes it ``easier on the eyes.''  For example,
% consider $\dx\left((a^2+x^2)^{\frac{1}{2}} \right).$ Our rules
% can handle differentiating a fractional exponent.  The problem is that
% what is in the base of that exponent is a complicated expression.

% As we said before, substitution in a powerful tool in nearly every
% branch of mathematics.  Utilizing it in differentiation as we have, it
% has a special name: The Chain Rule.  As such, it is more than just
% symbolic manipulation and provides us with a link between Leibniz'
% static geometric view of Calculus and Newton's dynamic physical
% approach.  We will explore what we mean by this, but for now,
% recognize that the Chain Rule really is substitution as you work on
% these practice problems. 

% See footnote\footnote{We need a couple of simple examples of how to derive
%   differential equations from a word problem. Consider the constant
%   subtangent equal to one (exponential), and $x^2+y^2=1 \imp
%   \dfdx{y}{x}=-\frac{x}{y}.$ Maybe one other?} }, suppose that $y=x^2.$ Then as we've seen $\dx y= 2x\dx x.$
% Now suppose that $x=t^3+1.$ Then $\dx x = 3t^2\dx t.$ Putting this all
% together gives:
% \begin{align*}
%   \dx y &= 2x\dx x\\
%        &= 2x(3t^2\dx t).\\
%   \intertext{At this point we have $\dx y$ in terms of both $x$ and $t.$
%              There is nothing wrong with this, \emph{per se,} but in
%              general we  prefer using only one variable. So we make
%              the substitution $x=t^3+1,$ giving:}
%   \dx y    &= 2(t^3+1)(3t^2\dx t).\\
% \end{align*}
% And finally, if we want $\dfdx{y}{t}$ we simply divide through by $\dx
% t$ giving:
% $$
% \dfdx{y}{t} = \underbrace{2(t^3+1)}_{\dfdx{y}{x}}(\underbrace{3t^2}_{\dfdx{x}{t}}).
% $$

% Notice that since $\dfdx{y}{x} =  2(t^3+1)$ (since $x= t^3+1$) and that
% $\dfdx{x}{t} = 3t^2$ this is just a particular example of the formula
% $$
% \dfdx{y}{t} =\dfdx{y}{x}\dfdx{x}{t}.
% $$ 
% Since this looks so much like the older ``Chain Rule'' for changing
% money\footnote{In fact, it really is the same rule if you look at it
%   in just the right way.} from one
% currency to another it has inherited the name ``The Chain Rule,'' and
% is usually stated as a separate theorem as follows.

% \begin{mytheorem}{The Chain Rule}
% \index{Differentiation Rules!Chain Rule}
%   \label{thm:chain-rule}
%   If $y$ depends on $x,$ and $x$ depends on $t$ then  
%   $$
%   \dfdx{y}{t} =  \dfdx{y}{x}\dfdx{x}{t}.
%   $$
% \end{mytheorem}
% \enddigress{Chain Rule}




% \section{The Chain Rule}
% \label{sec:chain-rule}
% Suppose for some reason we want to find the number of seconds in a
% typical $365$ day year\index{seconds per year}.  For our purposes we are not so interested in
% the actual number, which will turn out to be approximately
% $\pi\times10^7,$ as we are in the computational process involved. We
% start with the observation that the proportion 
% \[
% \left(60\
%   \frac{\text{seconds}}{\text{minute}}\right)\times\left(60\
%   \frac{\text{minutes}}{\text{hour}}\right)
% \] 
% yields \(3600\
% \frac{\text{seconds}}{\text{hour}}.\) Notice that the ``seconds'' that
% appear in the numerator of the first factor and the denominator of the
% second factor divide out\footnote{Sometimes this is called
%   ``cancelling'' but we will avoid this term as it is the source of
%   many errors.}, leaving only the proportion, seconds per hour. In the next
% step we multiply \(3600\ \frac{\text{seconds}}{\text{hour}}\) by \(24\
% \frac{\text{hours}}{\text{day}}\) to get the number of seconds per
% day. Continuing in this fashion and ``chaining'' the computations
% together we get:
% \[
% \left(60
%   \frac{\text{\small{}seconds}}{\text{\small{}minute}}\right)\times\left(60
%   \frac{\text{\small{}minutes}}{\text{\small{}hour}}\right)\times\left(24
%   \frac{\text{\small{}hours}}{\text{\small{}day}}\right)\times\left(365 \frac{\text{\small{}
%       days}}{\text{\small{}year}}\right)=31,536,000
% \frac{\text{\small{}seconds}}{\text{\small{}year}}.
% \]
% This kind of computation was called the Chain Rule for centuries
% before Calculus came along and it is easy to see why. We can get from
% one kind of proportion to another as long as we can find a  chain of
% proportions where all of intermediate numerators and denominators will
% divide out.
% \begin{wrapfigure}[]{R}{2.5in}
% %\vskip1.2cm{}
%   \includegraphics*[height=1.2in,width=2.5in]{../Figures/ChainRule}
% \caption{\textcolor{blue}{The Chain Rule: Notice that $y$ depends on $x$ which depends
%   on $t.$}}
% \label{fig:ChainRule}
% %\vskip1mm{}
% \end{wrapfigure}

% Now consider the situation depicted in figure~\ref{fig:ChainRule}
% where $y$ depends on $x$ and $x$ depends on $t.$ Suppose we know
% $\dfdx{y}{x},$ the rate of change of $y$ with respect to $x,$ and we
% know $\dfdx{x}{t},$ the rate of change of $x$ with respect to $t,$ but
% we need $\dfdx{y}{t},$ the rate of change of $y$ with respect to $t.$
% Clearly we need to compute
% \begin{equation}
% \dfdx{y}{x}\dfdx{x}{t}=\dfdx{y}{t}.
% \label{eq:ChainRule}
% \end{equation}\index{Chain Rule}
% Since this is a Chain Rule type of conversion, when Calculus was invented
% this particular computation was also called the Chain Rule. As time
% went on the older meaning dropped away so that today this particular
% operation with differentials is the only Chain Rule. 

% We formalize this as a theorem.

% \begin{mytheorem}{The Chain Rule}
%   If $y$ depends on $x,$ and $x$ depends on $t$ then the rate of
%   change of $y$ with respect to $t$ is given by the formula:
% \[
% \dfdx{y}{t} = \dfdx{y}{x}\cdot\dfdx{x}{t}.
% \]
% \end{mytheorem}

% There is more in our symbolism than is at first
% apparent. Consider the following very simple example.
% \begin{myexample}{}
%   \label{example$:cr1}
%$   Suppose that $y=3x$ and that $x=5t.$ Compute $\dfdx{y}{t}.$

% This is such a simple problem that we will have no problem computing
% $\dfdx{y}{t}$ directly. Since
% \begin{align*}
%   y&=3x\\
%    &=3(5t)\\
%    &=15t\\
% \intertext{we have}
%  \dx y &=15\dx t \text{ or}\\
% \dfdx{y}{t}&=15
% \end{align*}
% so the rate of change of $y$ with respect to $t$ is $15.$ 

% But computing $\dfdx{y}{t}$ directly like this hides most of what is
% interesting and all of what will ultimately be useful to us.  

% So let's do it again by the following, admittedly more cumbersome
% method.  Looking again at equation(~\ref{eq:ChainRule}) we see that we
% need $\dfdx{y}{x},$ which for this problem is $\dfdx{y}{x}=3.$  We
% will also need $\dfdx{x}{t}$ which for this problem is $\dfdx{x}{t}=5.$
% Thus
% \[
% \dfdx{y}{t} = \dfdx{y}{x}\cdot \dfdx{x}{t} = 3\cdot5 = 15.
% \]
% Naturally we get the same answer as before. Only the method of
% computation is different. 
% \end{myexample}

% Example~\ref{example:cr1} is a little too simple to display all of the
% nuances of the Chain Rule. The next example is slightly more complex.

% \begin{myexample}{}
%   \label{example:cr2}
% Now suppose $y=3x^2$ and again $x=5t.$ Let's compute $\dfdx{y}{t}.$

% First we need $\dfdx{y}{x}.$ To find this we compute 
% \[\dx y = d\left(3x^2\right).\] You may already see that 
% \[\dx y = 6x\dx x.\] If so, good for you! However, let's walk through
% the details of this calculation just to be sure. 

% Clearly the expression $3x^2$ is the product of $3$ and $x^2$ so by
% the Product Rule we have
% \[
% \dx y =3\dx (x^2)+x^2\dx (3).
% \]
% We saw on page~\ref{thm:ConstantDifferential}
% that the differential of a constant is zero so $\dx (3)=0.$ This leaves 
% \[\dx y = 3\dx (x^2). \]
% If you did Problem~\#\ref{prob:2} you already know that $\dx (x^2)=2x\dx
% x.$ If you haven't done Problem~\#\ref{prob:2},  the calculation is
% simple enough. Simply observe that $x^2$ is the product of $x$ and $x$
% and apply the Product Rule again. Either way we now have
% \[
% \dx y = 3(2x\dx x) = 6x\dx x
% \]
% or 
% \begin{equation}
%   \label{eq:CRExample2}
%   \dfdx{y}{x} = 6x.
% \end{equation}
% We also need 
% \begin{equation}
%   \label{eq:CRExample2-1}
%   \dfdx{x}{t} =5.
% \end{equation}
% Combining equations~\ref{eq:CRExample2} and~\ref{eq:CRExample2-1}
% gives
% \begin{align}
% \nonumber  \dfdx{y}{t} &= \dfdx{y}{x}\cdot\dfdx{x}{t}\\
% \nonumber  &= (6x)(5)\\
% \intertext{so that}
% \dfdx{y}{t} &= 30x.  \label{eq:CRExample2-3}
% \end{align}

% This seems like the solution to our problem until we think about it
% for a moment. Then we realize that the symbol $\dfdx{y}{t}$ means
% ``the rate of change of $y$\index{with respect to} \underline{with
%   respect to $t.$''} 

% It is this last phrase, ``with respect to $t$,'' which is
% troubling. When we say that $y$ depends on $t$ we are thinking of $t$
% as the independent and $y$ as the dependent variable so the right side
% of equation(~\ref{eq:CRExample2-3}) should be in terms of $t$ not $x.$

% This is easily handled. Since $x=5t$ we have 
% \[
% \dfdx{y}{t} = 30x = 30(5t) = 150t.
% \]
% \end{myexample}

% Again, the next example is slightly more complex than those previous.

% \begin{myexample}{}
%   \label{example:CR3}
% Suppose $y=4x^2$ and $x=3t^2.$ We wish to compute $\dfdx{y}{t}.$ As
% before we need both $\dfdx{y}{x}$ and $\dfdx{x}{t},$ since 
% \[
% \dfdx{y}{t} = \dfdx{y}{x}\cdot\dfdx{x}{t}.
% \]

% Thus
% \begin{align*}
%   \dx y &= \dx(4x^2)\\
%   &= 4\dx(x^2)\\
%   &= 4(2x)\dx x\\
%   &= 8x\dx x \\
% \intertext{or}
% \dfdx{y}{x} &= 8x.
% \end{align*}

% Similarly, 
% \begin{align*}
%   \dx x &= \dx(3t^3)\\
%   &= 3\dx(t^3)\\
%   % &= 3\dx(t^2\cdot t)\\
%   % &= 3\left[t^2\dx t+t\dx(t^2)\right]\\
%   % &= 3\left[t^2\dx t+t(2t)\dx t\right]\\
%   % &= 3\left[t^2\dx t+2t^2\right]\dx t\\
%   &= 3\left(3t^2\right)\dx t\\
% \intertext{so that}
% \dfdx{x}{t} &= 9t^2. \\
% \intertext{Thus}
% \dfdx{y}{t}&=\dfdx{y}{x}\cdot\dfdx{x}{t}\\
%   &= (8x)(9t^2).\\
% \intertext{As before the symbol $\dfdx{y}{t}$ indicates that we are
%   thinking of $t$ as the independent variable so we substitute
%   $x=3t^2$ to get}
% \dfdx{y}{t} &= 8(3t^2)(9t^2)\\
% \intertext{or}
% \dfdx{y}{t}&=216t^4.
% \end{align*}

% \end{myexample}
% \begin{ProblemSection}
%   \begin{embeddedproblem}{}
%     Verify the computations in Example~\ref{example:CR3} by direct
%     substitution.
%   \end{embeddedproblem}
% \end{ProblemSection}


%%% Local Variables: 
%%% mode: latex
%%% outline-minor-mode: t
%%% eval:(flyspell-mode)
%%% TeX-master: "DiffCalc"
%%% End: 
